
\begin{minipage}{0.38\textwidth}
\begin{center}
\includegraphics[scale=1,trim={0 1cm 0 4cm}, clip]{img/one_on_one/zonal_blocking/position}

{\it The track may be divided using cones. The blocker is strongly encouraged to suck back past the front cones into the jammer's space.} 
\end{center}
\end{minipage}
\hfill
\hspace{0.5cm}
\begin{minipage}{0.6\textwidth}

\subsection*{Drill: Zonal Catching}
\label{drill:one_on_one/zonal_catching}

% Nova Drill

Divide the track into thirds, with a middle, inside and outside lane. 
Place cones on these dividing lines marks, spaced about 50cm apart.

This drill requires two skaters, one acting as a jammer, one as a blocker, however this is a blocker focused drill. 
The jammer should sequentially move backwards around the cones, then attempt to pass the blocker in the next lane. 
The blocker should move between lanes, tracking and catching the jammer.

The goal of this drill is the footwork and the catch rather than the jammer attempting to pass the blocker.

The drill may be run with any form of catch.

\end{minipage}
